\documentclass[12pt]{article}
% Préambule commun à tous les documents extraits de « La Revue CAPES »
\usepackage[T1]{fontenc}
\usepackage[utf8]{inputenc}
\usepackage{lmodern}
\usepackage{amsmath,amssymb,amsthm,amsfonts,mathrsfs}
\usepackage{setspace}
\usepackage{listings}
\usepackage{pgfplots}
\pgfplotsset{compat=1.18}
\usepackage{longtable}
\usepackage{adjustbox}
\usepackage{lettrine}
\usepackage{tkz-tab}
\usepackage{changepage}
\usepackage{pdflscape}
\usepackage{graphicx}
\usepackage{tikz}
\usepackage{xcolor}
\usepackage[a4paper,top=2cm,bottom=2.5cm,left=2.5cm,right=2cm]{geometry}
\usepackage{fancyhdr}
\usepackage{draftwatermark}
\SetWatermarkText{La RevueCapes}
\SetWatermarkLightness{0.9}
\SetWatermarkScale{2}
\usepackage{hyperref}
\hypersetup{colorlinks=true,linkcolor=blue!50!black,urlcolor=blue!50!black}

\definecolor{revue}{RGB}{20,60,120}

\pagestyle{fancy}
\fancyhf{}
\renewcommand{\headrulewidth}{0.4pt}
\setlength{\headheight}{15pt}

% \entete{Matière}{Titre}{Type de document}
\newcommand{\entete}[3]{%
  \fancyhead[L]{\small\textsc{La Revue CAPES} -- #1}%
  \fancyhead[R]{\small #2}%
  \fancyfoot[C]{\thepage}%
  \thispagestyle{plain}%
  \begin{center}
    {\color{revue}\rule{\textwidth}{1.2pt}}\\[4pt]
    {\small\textsc{Concours directs CAP-CEG / CAPES -- Burkina Faso}}\\[6pt]
    {\LARGE\bfseries\color{revue} #2}\\[6pt]
    {\large\itshape #1 -- #3}\\[2pt]
    {\color{revue}\rule{\textwidth}{1.2pt}}
  \end{center}
  \vspace{0.5cm}%
}

% Encadré pour les remarques ajoutées dans les corrigés
\newcommand{\remarque}[1]{\par\medskip\noindent\fcolorbox{revue}{revue!6}{\parbox{\dimexpr\linewidth-2\fboxsep-2\fboxrule}{\textbf{\color{revue}Remarque.} #1}}\par\medskip}


\begin{document}
\entete{Mathématiques}{Ancien sujet CAPES -- Session 2017}{Énoncé}
\begin{spacing}{1.5}

\textbf{Exercice 1 (6pts)}


On considère la matrice carrée d'ordre 3 définie par  : 

\[M=\begin{pmatrix}
11 & -5 & 5 \\ 
-5 & 3 & -3 \\ 
5 & -3 & 3
\end{pmatrix}\]

1. Déterminer le polynôme caractéristique de $M$, ainsi que son spectre.

2. Quels sont les sous espaces propres associés aux différentes valeurs propres?

3. $M$ est-elle diagonalisable ? Si oui, la diagonaliser. On précisera la base $B$ et les matrices $D$ et $P$ de cette base.

4. En déduire $M^{n}$ $\forall n\in \mathbb{N}$.


\medskip
\textbf{Exercice 2 (8pts)}

On considère l'équation différentielle :

\begin{center}
$2y''-y'-y=-7-3x+(12x+5)e^x$.
\end{center}

1. Résoudre l'équation homogène associée à $(E)$. On notera $y_H$ sa solution.

2. Trouver une solution particulière de l'équation suivante sous le forme $y_1=mx+p$ avec $m,p\in \mathbb{R}$ :

$2y''-y'-y=-7-3x$.

3. Trouver une solution particulière de l'équation suivante sous le forme 

$y_2=x(ax+b)e^x$ avec $a,b\in \mathbb{R}$

$2y''-y'-y=(12x+5)e^x$.

4. En déduire une solution particulière $y_p$ de l'équation $(E)$.

5. Écrire la solution générale $y_G$ de l'équation différentielle $(E)$.

6. Déterminer la solution $f$ de $(E)$ telle que la tangente à la courbe représentative de $f$ au point $A(0,2)$ soit parallèle à la droite $(D)$ d'équation $y-9x+9=0$.

\medskip
\textbf{Exercice 3 (6pts)}

Soit $f$ une fonction définie sur $[-1,+\infty[$ par :

\[f(x)=\frac{2}{\pi}\int_{0}^{\frac{\pi}{2}}\sqrt{1+x\cos^2t}dt\].

1. Préciser $f(-1)$ et $f(0)$.

2. Établir pour tout $x>-1$, la relation $f(x)=\sqrt{1+x}f(\frac{-x}{1+x})$.

3. Démontrer que pour $-1\leq x_0\leq x$, on a l'encadrement :

$0\leq f(x)-f(x_0)\leq \frac{x-x_0}{4}$.

4. Démontrer que pour $-1<x<0$, on a l'encadrement 

$0\leq f(x)-f(-1)\leq \frac{1}{2}\sqrt{1+x}$.

5. Étudier la continuité de $f$.



\medskip

\end{spacing}
\end{document}
